% Paper 2 	extperiodcentered{}  §5 Discussion
% Status: draft 	extperiodcentered{}  2026-05-05

\section{Discussion}
\label{sec:discussion}

\subsection{When does each component help?}

The 5 components in v0.8 are independent and can be ablated in
isolation. Their effectiveness varies by question type:

\begin{itemize}
\item \textbf{TOP\_K 10\(\to\)15}: marginal gain (\(+0.5\) pt) across all
  types. Diminishing returns above 15 in our pilot (\(+0.0\) at K=20,
  \(-0.5\) at K=30 due to context contamination).

\item \textbf{ssa context expansion (2400\(\to\)3500 chars)}: \(+2\) pts on
  ssa, neutral elsewhere. ssa questions explicitly require fine-grained
  recall of one specific session; the truncation removed answer-relevant
  detail.

\item \textbf{multi-session decompose}: \(+8\) pts. The LLM reliably
  miscounts when given 5+ sessions in flat form; explicit per-session
  enumeration in the prompt fixes this.

\item \textbf{knowledge-update timestamp}: \(+2\)--\(3\) pts. Helps but
  not dramatic; the LLM tends to honor the most-recent statement
  by default. Explicit instruction tightens that prior.

\item \textbf{Multi-angle query rewriting (ssu)}: \(+27\) pts. The
  largest single gain. ssu questions are referent-light (``what dish
  cannot the user eat?'') and dense retrieval misses without
  rephrasing. We expect this gain to transfer to other benchmarks
  where queries are colloquial.
\end{itemize}

\subsection{The temporal-reasoning bottleneck}

Temporal-reasoning is the only question type where v0.8 does not
significantly improve over baseline (\(+0.7\) pt). Inspection of failure
cases shows the LLM correctly retrieves the relevant sessions but
miscomputes time arithmetic (e.g., ``how many days between sessions
3 and 7?''). We hypothesize this is a fundamental LLM weakness rather
than a retrieval issue and is the right place for future targeted work
(structured time extraction, calendar tools).

\subsection{Trajectory observations during the full-500 run}

Plotting cumulative accuracy by question index reveals a V-shape:

\begin{itemize}
\item \texttt{ssu} batch (Q1--70): peak 60.8\% (\(+30\) pts vs baseline)
\item \texttt{multi-session} batch (Q70--204): drops to 55.5\%
\item \texttt{ssp} batch (Q204--234): 56.2\% (default prompt)
\item \texttt{temporal-reasoning} batch (Q234--367): trough at 48.2\%
\item \texttt{knowledge-update} batch (Q367--444): rebound to 52.8\%
\item \texttt{ssa} batch (Q444--500): peak \(\sim\)84\%
\end{itemize}

The trough at temporal-reasoning is a per-type difficulty effect, not a
model degradation; the V-shape recurs in every model we tested.

\subsection{Anchor-based drift detection (orthogonal contribution)}

Beyond the LongMemEval pipeline, Compass v0.8 includes an anchor-based
drift detector that runs at hook time on every user prompt
(\S~A~appendix). The detector embeds the prompt and computes cosine
distance to 25 positive anchors (intended behavior) and 35 negative
anchors (drift exemplars). This is orthogonal to LongMemEval performance
but is the second contribution claim of this paper:

\begin{itemize}
\item AUC = 0.92 on a held-out 200-prompt test set
\item Latency = \(<\)50\,ms p95 (BGE-m3 daemon, batched)
\item Per-domain anchor packs (finance, legal, medical, vc, zenmind)
  ship pre-built; user-domain anchors auto-train from
  feedback signal
\end{itemize}

This makes Compass the first open-source memory system that combines
recall with drift detection at hook latency, in contrast to
narrative-summary memories (claude-mem) and graph memories (Zep)
that focus only on retrieval.

\subsection{Cost economics for industry adoption}

For a Chinese-region production deployment (Tencent T4 spot + Volc Ark
coding plan), running compass v0.8 in 24/7 hook mode costs:

\begin{itemize}
\item GPU: \cny{}300/month (1 T4 spot, \(\sim\)\$45)
\item LLM API: \cny{}50--500/month per active user (varies with
  session frequency)
\item bge-m3 inference: 0 marginal cost (local, daemon)
\end{itemize}

For the same workload using only commercial APIs (GPT-4o + Claude
Sonnet), costs would be \(\geq\)\(20\times\) higher. We argue that for
the Chinese market specifically, a fully-local-bge + Volc-Ark-LLM
stack achieves SOTA accuracy at price-points compatible with
\(\geq\)100,000 monthly active users on a small SaaS budget.

\subsection{Limitations and threats to validity}

(see Section~\ref{sec:limitations})
